\title{Fewer is More: Boosting LLM Reasoning with Reinforced Context Pruning}

\begin{abstract}
Large Language Models (LLMs) have shown impressive capabilities, yet they still struggle with math reasoning. In this work, we propose \textbf{CoT-Influx}, a novel approach that pushes the boundary of few-shot Chain-of-Thoughts (CoT) learning to improve LLM mathematical reasoning. Motivated by the observation that adding more concise CoT examples in the prompt can improve LLM reasoning performance, {CoT-Influx} employs a coarse-to-fine pruner to maximize the input of effective and concise CoT examples. The pruner first  selects as many crucial CoT examples as possible and then prunes unimportant tokens to fit the context window. A math reasoning dataset with diverse difficulty levels and reasoning steps is used to train the pruner, along with a math-specialized reinforcement learning approach. As a result, by enabling more CoT examples with double the context window size in tokens, {CoT-Influx} significantly outperforms various prompting baselines  across various LLMs (LLaMA2-7B, 13B, 70B) and 5 math datasets, achieving up to 4.55\% absolute improvements. Remarkably, without any fine-tuning, LLaMA2-70B with {CoT-Influx} surpasses GPT-3.5 and a wide range of larger LLMs (PaLM, Minerva 540B, etc.) on the GSM8K. {CoT-Influx} serves as a plug-and-play module for LLMs and is compatible with most existing reasoning prompting techniques, such as self-consistency and self-verification.
\end{abstract}

\section{Introduction}

Large Language Models (LLMs) have demonstrated remarkable capabilities across a range of tasks~\cite{gpt3,gpt4}. However, it remains a significant challenge to improve LLM performance on reasoning tasks, especially for smaller LLMs like LLaMA~\cite{llama} on math reasoning. 

While existing efforts focus on optimizing Chain-of-Thought (CoT) prompts~\cite{cot, selfconsistency,treeofthought} and fine-tuning LLMs~\cite{luo2023wizardmath} under the zero-shot setting, the potential of few-shot learning in improving LLM reasoning has not been fully explored. Inspired by the human reasoning process, we propose the hypothesis: \textit{if LLMs are exposed to more step-by-step problem-solving examples (i.e., CoTs) before answering questions, it could potentially improve LLMs reasoning capability to generate a correct solution}.  This leads to our question: \textit{what's the boundary of LLM reasoning capability achievable through inputting more CoT examples?}

However, we face two major obstacles. First, the limited token length of LLMs' context window restricts the number of few-shot examples. Extending the context window is one solution, but it requires expensive fine-tuning and increases inference overhead~\cite{pi,peng2023yarn}. While prompt compression~\cite{li2023compressing,llmlingua} is another approach, it underperforms in math reasoning. Tokens like numerical and format ones, though identified redundant, are crucial for few-shot math problem solving. 

Second,  {it's challenging to select helpful CoT examples}. 
Section~\ref{sec:pilot_study} reveals that random choices can even harm reasoning performance. Existing retrieval-based methods~\cite{topk,scarlatos2023reticl} are not tailored  for math reasoning, making them suboptimal. These retrieved examples are model-agnostic, while we found that different LLMs  favor CoT examples of varying characteristics (e.g., diverse difficulty levels).

In this work, we propose \textbf{CoT-Influx}, which addresses all the above challenges and pushes the boundaries of utilizing few-shot learning to improve  LLM math reasoning capability. {CoT-Influx} is motivated by the observation that \textit{current LLM context window has not been fully utilized due to redundancy at both the example and token levels in natural language input}. As such, these redundant inputs can be pruned to free up space for more informative context. The central idea of {CoT-Influx} is to input long lengthy CoT examples, select the crucial examples for the target LLM, and then prune redundant tokens to fit within the original LLM context window. As a result, by inputting much more helpful CoT examples, each composed solely of informative tokens and with a shorter length, we greatly improve LLM ability to solve math problems. Moreover, as all these inputs remain within the context window, we do not increase any inference overhead. This stands in stark contrast to other methods~\cite{hao2022structured,pi}.

{CoT-Influx} treats the target LLM as a black box, and serves as a plug-and-play module for LLMs as shown in Fig.~\ref{fig:overview}. The key module is a coarse-to-fine pruner involving two steps: (i) a shot pruner first selects the most helpful CoT examples from a large batch of shots,  and (ii) a token pruner then removes unimportant tokens from these selected CoT examples. To effectively train the pruner module tailored for math reasoning, {CoT-Influx} is built upon the following novel techniques. 

First,  {CoT-Influx} requires a CoT dataset for training and inference. Existing CoT examples,  heavily reliant on costly human engineering, often struggle with diversity and quality.  To address this, we employ GPT-4~\cite{gpt4} and Evol-Instruct~\cite{xu2023wizardlm} to create a math reasoning dataset, called \textit{MRD$^3$}.
With problems of varying difficulty and reasoning steps, MRD$^3$ enables  {CoT-Influx} to generalize across a wide range of math problems. 

Second, training the pruner presents two challenges: (1) since we identify discrete tokens before the LLM tokenizer, the loss gradient cannot be backpropagated through the tokenizer to update the pruner; (2) The high difficulty of many math problems, which consistently yield incorrect answers regardless of the quality of compressed few-shot examples, poses a challenge to the effective training of the pruner. To this end, we introduce a novel training approach with reinforcement learning to mitigate the gradient issue. We design a reward function to measure the LLM loss, few-shot math reasoning effectiveness, and token length constraints. Then, we design a difficulty-aware dataloader filtering appropriate problems and introduce two techniques to stabilize the RL training. 

Extensive experiments on various LLMs and five math datasets demonstrate the effectiveness of {CoT-Influx}. {CoT-Influx} significantly boosts LLM reasoning capability,  achieving 1.36\%-14.09\% absolute improvements over SOTA baselines, and establishes a new prompting-based benchmark in math reasoning accuracy without any fine-tuning or additional inference costs. Remarkably, LLaMA2-70B with {CoT-Influx} outperforms a broad range of larger LLMs and surpasses GPT-3.5 by 2.5\% on GSM8K. Moreover, {CoT-Influx}  excels over retrieval and prompt compression baselines in example selection and identifying crucial tokens.

\section{Related Works}

\textbf{LLMs for Math Reasoning}. Drawing from the Chain-of-Thought (CoT)~\cite{cot}, recent research has greatly improved the reasoning capabilities of LLMs by providing step-by-step reasoning paths. The main efforts are twofold: enhancing CoT prompts, such as Program-of-Thoughts~\cite{pot}, Tree-of-Thoughts~\cite{treeofthought}, and Everything-of-Thoughts~\cite{ding2023everything}, and innovating CoT-based training data for fine-tuning LLMs like WizardMath~\cite{luo2023wizardmath}. 

However, most works focus on the zero-shot setting with only task instruction or CoT prompts,  leaving the potential of few-shot CoT largely untapped. We explore leveraging few-shot CoT learning to improve LLMs' math reasoning capabilities.

\textbf{Prompt Compression}. To address the challenge of limited few-shot examples due to restricted context window length, one related work involves prompt compression.  Key approaches  include:  (1) token pruning~\cite{ltp,top}; (2) soft prompt compression methods~\cite{wingate2022prompt,mu2023learning,chevalier2023adapting,ge2023context}; and (3) information-entropy-based approaches~\cite{li2023compressing,llmlingua}.

 However, they do not effectively solve our problem for two reasons. First, they prune tokens based on designed metrics, often failing to remove  redundancy of the entire CoT examples. Second, some tokens such as numerical and format tokens, although redundant, are crucial for math reasoning. 

\textbf{Prompt Retrieval} optimizes task performance by selecting high-quality few-shot examples using either heuristics or a supervised retriever model. Heuristic methods, such as the widely used TopK retrieval~\cite{topk, gao2021making}, BM25~\cite{robertson2009probabilistic}, VoteK~\cite{hongjin2022selective}, and entropy~\cite{lu2022fantastically}, select examples based on semantic similarity. Recently, supervised-based methods like  EPR~\cite{rubin2021learning}, LLM-R~\cite{wang2023learning}, and IDS~\cite{qin2023context} have been proposed, which train a retrieval model to learn better example selection.

However, these methods are sub-optimal for math reasoning, as they retrieve model-agnostic examples. In contrast, LLMs with different capabilities favor CoT examples of varying complexities. Moreover, they don't account for token redundancy, which restricts the number of retrieved examples. 

\section{Pilot Study}\label{sec:pilot_study}

This section presents our key observations of few-shot learning in improving LLMs math reasoning, upon which the {CoT-Influx} design is based. Note that experiments are done with our proposed CoT dataset, {MRD$^3$}, as introduced in Sec.~\ref{sec:dataset}.

\textbf{Observation 1}: \textit{LLMs can improve reasoning with more helpful CoT examples, but the current context window restricts the number of  CoT examples}.

A standard practice for evaluating LLMs' math reasoning capability is the use of 8-shot manually-designed CoTs~\cite{cot}. We increase the number of CoT shots to see if reasoning accuracy improves. To avoid poor-quality examples, we use the TopK method~\cite{topk} to select the $k$ most relevant CoT examples for each question. Given LLaMA2's context window limit of 4096 tokens, we could only input up to 20 CoT examples\footnote{The input token length is less than the context window token limit, as the answer generation also shares this limit. }. As Fig.~\ref{fig:cotnumber} shows, increasing CoT examples improves LLaMA2-7B's reasoning accuracy on the GSM8K dataset, significantly outperforming the standard 8-shot setting.  However, the limited LLM context window hinders the full potential of few-shot CoT learning for improving math reasoning. For instance, even with 20 CoTs not hitting the token limit, accuracy drops as the large input context limits the LLM's response space.

\textbf{Observation 2}: \textit{CoT example selection is crucial for math reasoning. Simply adding  CoT examples randomly doesn't boost performance}.

The prior study suggests that more CoT examples can improve LLM reasoning performance. However, the quality of CoT examples is crucial to the final performance. As shown in Table~\ref{tbl:cotquality}, even with up to 16 CoT shots, random selection underperforms the standard 8-shot setting, which is manually curated for quality. 

\textbf{Observation 3}: \textit{A CoT example contains redundant tokens for math reasoning, which can be pruned to free up space for more informative content.} 

Observation 2 indicates that few-shot CoT examples contain useless or even harmful examples that can be pruned. We further observe that a CoT example often has redundant tokens. For instance, the blue tokens in Fig.~\ref{fig:compressed_example} can be removed without affecting LLM performance. However, identifying redundant tokens for math reasoning poses a challenge. Simply using existing prompt compression methods~\cite{llmlingua,li2023compressing} leads to a significant performance decline. Fig.~\ref{fig:compressed_example} shows a compressed example using LLMLingua~\cite{llmlingua}. Some numerical and format tokens (colored in red), while identified as redundant, are crucial for LLM to comprehend the context for solving a math problem. 

\section{{CoT-Influx} Methodology}

Motivated by our observations, this section introduces {CoT-Influx}, which maximizes CoT examples within the LLM context window by identifying the most important CoT examples and tokens from long lengthy input contexts. 

\subsection{CoT Dataset Collection}

\label{sec:dataset}
We start by collecting a high-quality math reasoning dataset, comprising diverse CoT examples with varying steps and difficulties. We merge the training set of GSM8K~\cite{gsm8k}, MAWPS, MAWPS-single~\cite{mawps}, and 1000 random examples from AQuA~\cite{ling2017program} to create an initial dataset of 9.7K question-answer pairs. Then, we prompt GPT-4 to generate formatted CoT reasoning steps. Notably, it's crucial to maintain a consistent format for each example in few-shot learning.  Our dataset also assigns a difficulty score from 1 to 10 for each question, based on GPT-4's evaluation, where 1 signifies the easiest questions and 10 is the most difficult. 

We observe that most questions in this initial dataset score between 2-4. To improve difficulty diversity, we use GPT-4 to mutate questions, generating corresponding CoTs with varied difficulty levels. We apply 5 mutation schemes, three to increase reasoning difficulty and two to simple questions. The final dataset is referred to as \textit{M}ath \textit{R}easoning \textit{D}ataset with \textit{D}iverse \textit{D}ifficulty (\textit{MRD$^3$}). 

\subsection{Problem Formulation}
\label{sec:problem_formulation}

Let $\mathcal{D}$ denote the  CoT dataset (i.e., the MRD$^3$), and $\hat{\mathcal{D}}=\{x^\text{cot}_i\}_{i=1}^k$ be a subset of $k$ CoT examples, each composed of a question, reasoning steps, and an answer. The total number of tokens in these $k$ CoT examples far exceeds the LLM context window length limit of $T$.    {CoT-Influx} is designed to perform a two-step pruning process:

\begin{equation}
\resizebox{0.43\textwidth}{!}{$
\hat{\mathcal{D}}=\{x^{\text{cot}}_i\}_{i=1}^{k} \xrightarrow{\text{Shot Pruner}} \{ x^{\text{cot}}_j\}_{j=1}^{k'} \xrightarrow{\text{Token Pruner}} \{ \hat{x}^{\text{cot}}_{j}\}_{j=1}^{k'}  $}
\end{equation}

Initially, non-useful CoT examples are pruned from $\hat{\mathcal{D}}$, resulting in a reduced set of $k'$ examples. Then, for each retained CoT example $x^\text{cot}$, redundant tokens are pruned, yielding a shorter example,  $\hat{x^\text{cot}}$.

Let $x^\text{question}$ denote the question that LLM is tasked to solve. For final input $x^\text{input}$, we concatenate all tokens from $\{ \hat{x}^{\text{cot}}_{j}\}_{j=1}^{k'}$ and place them before $x^\text{question}$. 
Our goal is to optimize the input $x^\text{input}$, so that LLM can correctly answer the question under $x^\text{input}$.

Meanwhile, the token count of $x^\text{input}$,  $t$, must be less than the LLM context window limit $T$. Formally, we optimize the following:

\begin{equation} \label{eq:target}
 	\resizebox{0.43\textwidth}{!}{$

 	\begin{aligned}
 			\min_{ \hat{\mathcal{D}}\subseteq \mathcal{D}} & L_{\text{LLM}}\left(x^{\text{input}}\right), \,
 			\max_{ \hat{\mathcal{D}}\subseteq \mathcal{D}} R_{\text{Acc}}\left(y_{\text{LLM}}\left(x^{\text{input}}\right), y^{\text{answer}}\right),
 			\\ s.t. \quad  &t\left(x^\text{input}\right)=\sum_{1}^{k'} |\hat{x}^{\text{cot}}|+|x^\text{question}|\leq T
 		\end{aligned}
 $}
 \end{equation}

where $L_\text{LLM}$ is LLM  loss, and $R_\text{Acc}$ evaluates if LLM's answer $y_{\text{LLM}}(x^{\text{input}})$ matches the correct answer $y^\text{answer}$, this will be elaborated in Sec.~\ref{sec:training}. 

\textbf{Overview}. Fig.~\ref{fig:overview} illustrates our approach. The core component is a lightweight,  plug-and-play module (Sec.~\ref{sec:prunerdesign}),  which consists of a small text embedding extractor and a coarse-to-fine pruner. 

To train the pruner, we face the challenge of gradient backpropagation when pruning discrete tokens outside the LLM. The  LLM loss gradient cannot be backpropagated through the tokenizer. To address this, we design a multi-objective reward function and use reinforcement learning for effective training (Sec.~\ref{sec:training}). The overall training process is outlined in Algorithm~\ref{alg:cot-max}.

\subsection{Coarse-to-fine Pruner Design}
\label{sec:prunerdesign}

\textbf{Text embedding extractor}. As {CoT-Influx} serves as an external module, we need to extract text embedding as prediction features. However, it's non-trivial to extract features for long inputs beyond the LLM context window. To address this, we use a small encoder model, BERT-Large~\cite{bert},  to extract sentence-level (i.e., a CoT example) embedding instead of extracting token embedding from the entire long context. For a batch of $k$ CoT examples, each example is padded to $N$=512 tokens. BERT then inferences these examples to obtain the final layer of text embedding, denoted as $H_{s_\text{shot}}\in\mathbb{R}^{k\times N\times D_{BERT}}$, where $D_{BERT}$ is BERT's hidden dimension size. 

\textbf{State}. As shown in Fig.~\ref{fig:overview}, we define state $s_\text{shot}\in\mathbb{R}^{k\times N}$ for the first shot pruner, representing the input batch of $k$ CoT examples $\in\hat{\mathcal{D}}$. For the second token pruner, we define state $s_\text{token}\in\mathbb{R}^{k'\times N}$, which represents all remaining tokens after the shot pruner. $k'$ is the number of retained examples.

\textbf{Action}. Let $a_\text{shot}$ and $a_\text{token}$ denote the actions predicted by the shot and token pruner, respectively. The action space is defined as  \{0, 1\}, where 1 signifies retention and 0 indicates pruning.  $a_\text{shot}$ determines whether each CoT example should be pruned, while $a_\text{token}$ predicts the pruning of each token in the retained CoT examples.

\textbf{Two-stage policy network}. The pruner module is a two-stage policy network, each stage is a two-layer feed-forward network (MLP) with GELU activation. This module outputs a continuous categorical distribution $\pi$,  used for sampling discrete actions (i.e., \{0, 1\}).  Let $\theta$ denote the MLP's trainable parameters and $\sigma(\cdot)$ the sigmoid function. Based on the current states $\{s_\text{shot}, s_\text{shot}\}$ and the hidden states $\{H_{s_\text{shot}}, H_{s_\text{token}}\}$, the policy network sequentially make two action predictions as follows:

\begin{equation} \label{eq:shot_policy}
	\small
\displaystyle \pi(a_\text{shot}|s_\text{shot};\theta) = \sigma\left( \text{MLP} \left( H_{s_\text{shot}}\right)\right)
\end{equation}

\begin{equation} \label{eq:token_policy}
	\small 
	\displaystyle
\pi(a_\text{token}|s_\text{token};\theta) = \sigma\left(\text{MLP}\left(H_{s_\text{token}}\right)\right),
\end{equation}

where $a_\text{shot}$ and $a_\text{token}$ are the predicted actions, sequentially predicting whether to prune each of the $k$ CoT examples and each token within the retained examples, respectively. We predict the discrete action by sampling from the categorical distribution $\pi$ with a softmax function.

\begin{algorithm}[t]
	\small
	\caption{Pruner Training and Inference}
	\label{alg:cot-max}
	\textbf{Input:} target LLM, dataset $\mathcal{D}$,   number of CoTs $k$, token limit $T$, manual few-shot cot $x^{\text{few-shot}}$, repeat $t_{\text{repeat}}$ \\

	\begin{algorithmic}[1]

		\STATE $\blacktriangleright$ \textbf{Phase 1: MRD$^3$ preperation}
		\STATE Perform evolution and difficulty evaluation to get $\mathcal{D}$;
		\STATE Use the difficulty filter and split $\mathcal{D}$ into questions set $\mathcal{D}_{\text{question}}$ and CoT candidates set $\mathcal{D}_{\text{cot}}$
		\STATE $\blacktriangleright$ \textbf{Phase 2: Training the two-stage pruner (1 epoch)}

		\FOR{$(x^{\text{question}},  y^{\text{answer}})$ in $\mathcal{D}_{\text{question}}$}

		\STATE Retrieve Top-$k$ examples $\hat{\mathcal{D}}=\{x^{\text{cot}}\}^k_{i=1}$ from $\mathcal{D}_{\text{cot}}$

		\STATE $H_{s_{\text{shot}}}$ = BERT$\left(\{x^{\text{cot}}\}^k_{i=1}\right)$ 

		\FOR{$j$=$1$ to $t_{\text{repeat}}$}

		\STATE Get $\pi\left(a_\text{shot}|s_\text{shot};\theta\right)$ with 
		Eq.~\ref{eq:shot_policy}, sample $a_\text{shot}$

		\STATE $\{x^{\text{cot}}\}^k_{i=1} \xrightarrow{}\{x^{\text{cot}}\}^{k'}_{i=1}$
		\STATE $H_{s_{\text{token}}}$ = BERT$\left(\{x^{\text{cot}}\}^{k'}_{i=1}\right)$

		\STATE Get $\pi\left(a_\text{token}|s_\text{token};\theta\right)$ with 
		Eq.~\ref{eq:token_policy}, sample $a_\text{token}$

		\STATE $\{x^{\text{cot}}\}^{k'}_{i=1}\xrightarrow{}\{\hat x^{\text{cot}}\}^{k'}_{i=1}$

		\STATE  $x^{\text{input}} = \left(\{\hat x^{\text{cot}}\}^{k'}_{i=1}, x^{\text{few-shot}}, x^{\text{question}}\right)$

		\STATE Output ${\text{LLM}}({x}^{\text{input}})$; Compute $R$ with Eq.~\ref{eq:total_reward}

		\ENDFOR

		\STATE Compute policy gradient using Eq.~\ref{eq:gradient}, update $\theta$ 
		\ENDFOR 

		\STATE $\blacktriangleright$ \textbf{Phase 3: LLM reasoning with pruner and MRD$^3$}

		\STATE Retrieve Top-$k$ shots $\{x^{\text{cot}}_{q}\}^k\in\mathcal{D}$ for  target question $q$

		\STATE Do pruning: $\{x^{\text{cot}}_{ q}\}^k\xrightarrow{\theta}\{\hat x^{\text{cot}}_{ q}\}^{k'}$, reconstruct $\{\hat x^{\text{cot}}_{q}\}^{k'}$ 

		\STATE $x^{\text{input}}_{q} = \left(\{\hat x^{\text{cot}}_{q}\}^{k'}, x^{\text{few-shot}}, x^\text{question}_q\right)$

		\STATE Get LLM reasoning output $ {\text{LLM}}({x}^{\text{input}}_{ q})$

	\end{algorithmic}
\end{algorithm}

\subsection{End-to-end RL Optimization}
\label{sec:training}
\textbf{Multi-objective Reward}. Our objective in Eq.~\ref{eq:target} is to train the pruner module to identify the most crucial CoT examples and useful tokens for math problem solving, while keeping the final tokens within the original LLM context window. To achieve this, we design a multi-objective reward.

Let $x^\text{input}$ be the final input to  LLM, which includes the retained CoT tokens from the policy network and the target question. $t$ represents the token count of $x^\text{input}$, and $T$ is the token count limit. The reward $R$ is defined as follows:

\begin{equation} \label{eq:total_reward}
	\small
	\displaystyle	R\left( x^\text{input}\right) = (\frac{1}{1+L_\text{LLM}}+R_{\text{Acc}}) \times \left[\frac{t}{T}\right]^w
\end{equation}

where the first term evaluates the effectiveness of inputted CoT tokens, and the second term ensures they are within the LLM context window. $L_\text{LLM}$ is LLM's prediction loss under $x^\text{input}$, $R_\text{Acc}$ evaluates the reasoning accuracy (to be discussed later). $w$ is a hyperparameter that penalizes the token count $t$ for being too short (i.e., $w<0$) or exceeding (i.e.,$w>0$ ) the token limit $T$. 

In addition to $L_\text{LLM}$, we introduce $R_\text{Acc}$ to evaluate math reasoning accuracy. This is because $L_\text{LLM}$,  the average loss of generated tokens, doesn't fully reflect LLM's ability to generate correct answers. Specifically, $R_\text{Acc}$ is set to 1 for a correct answer and 0 for an incorrect one. Notably, we found that if the format or crucial tokens are pruned, LLM struggles to interpret the input context correctly, leading to irrelevant answers for math problem solving.   In such cases, we penalize $R_\text{Acc}$ with a value of -0.1.

\textbf{Optimization with REINFORCE}. We employ reinforcement learning to maximize the reward and train the two-stage policy network. According to REINFORCE~\cite{reinforce}, the network parameters are updated by the gradients:

\begin{equation} \label{eq:gradient}
	\small
	\resizebox{0.35\textwidth}{!}{$ R \cdot \nabla_\theta \text{log} \pi(a_\text{shot}|s_\text{shot})\pi(a_\text{token}|s_\text{token})$} 
\end{equation}

Notably, as shown in Fig.~\ref{fig:overview}, only the parameters of the policy network require training. The embedding extractor and LLM are frozen, thus, the overall training overhead is lightweight.

\textbf{Difficulty-aware data filter}. Existing LLMs, particularly smaller ones, underperform in math reasoning. If the question is too challenging for LLMs, the answer will always be incorrect, regardless of the quality of compressed few-shot CoT examples, making it challenging to effectively train our pruner module. To address it, we use a difficulty filter to sample a math question set $\mathcal{D}_{\text{question}}$ from $\mathcal{D}$, which includes only easy questions with a difficulty score below a threshold $d_{\text{thr}}$.  During training, each question in $\mathcal{D}_{\text{question}}$ samples a batch of $k$ CoT examples from $\mathcal{D}_{\text{cot}}$, where $\mathcal{D}_{\text{cot}}$ is the CoT candidate set sampled from $\mathcal{D}$. 

\textbf{Stabilize the training}. Another challenge is that pruning CoT and tokens during training introduces instability, making it difficult for effective training. 

First, despite the optimization of  question set $\mathcal{D}_{\text{question}}$ through the filter, LLM performance for a randomly sampled question under different few-shot prompts can still be unpredictable. This unpredictability, where a question might yield correct results under low-quality pruned prompts and a complex question might fail under carefully pruned prompts, can affect the pruner's training effectiveness. To address this, we continuously repeat a sampled question multiple times, $t_\text{repeat}$, each time with a different pruned few-shot prompt from the pruner.  Moreover, we use exponential moving average (EMA) to smooth reward $R_\text{Acc}$ in Eq.~\ref{eq:total_reward}. 

Second, during the early training, our pruner module makes random decisions, leading to arbitrary removal of CoT examples and tokens. These randomly pruned few-shot prompts can cause instability in RL training. Empirically, we append the manually-designed 8-shot CoTs~\cite{cot} to the pruned prompts. This ensures a good lower bound and stabilizes the training. 

\section{Evaluation}

\textbf{Models, datasets and metric}. We evaluate  {CoT-Influx} on LLaMA2-7B, LLaMA2-13B, and LLaMA2-70B. The mathematical datasets  for evaluation  include GSM8K~\cite{gsm8k}, AddSub~\cite{addsub}, Multiarith~\cite{multiarith}, Svamp~\cite{svamp}, and Singleeq~\cite{singleeq}.  For evaluation metric, we report Exact Match (EM) accuracy of the predicted answers. 

\textbf{Baselines} We set  three  baselines for comparison:

\begin{itemize}[itemsep=2pt,topsep=0pt,parsep=0pt]
    \item \textit{CoT and few-shot CoT prompting}:  We compare with widely-used prompts for LLM reasoning, including zero-shot, zero-shot-CoT~\cite{kojima2022large}, and the standard few-shot-CoT~\cite{cot} with 8 shots. 
    \item \textit{Prompt retrieval}: we also compare with retrieval baselines, specifically using random, TopK~\cite{topk}, and BM25~\cite{robertson2009probabilistic} methods. We select as many CoT examples as possible using each method, without exceeding LLM context window. Random retrieval is   to reflect the average quality of our CoT dataset.  
    \item \textit{Prompt compression}: To evaluate the effectiveness of identifying crucial tokens, we compare the resulting compressed prompts from the same batch of CoT shots with state-of-the-art prompt compression baselines: Selective Context~\cite{li2023compressing}, LLMLingua~\cite{llmlingua}, and compression through GPT-4. 
\end{itemize}

\subsection{Main Results} \label{sec:main_results}

\textbf{Effectiveness of enabling more few-shot CoTs}. We first evaluate how far the boundary of few-shot learning can be pushed using {CoT-Influx}. For comparison, we set up two baselines: \textit{(i)} Few-shot CoT, using 8 manual-designed CoT shots as the default LLM evaluation setting on GSM8K. 
\textit{(ii)} TopK retrieves 20 CoT shots from our dataset, denoting the max shot number within LLaMA2 context window.  

For {CoT-Influx}, we test LLaMA2 7B and 13B on GSM8K, adjusting  the number of CoT shots from 16 to 64 examples, which corresponds to 0.7$\times$ to 2.8$\times$ the token count of LLaMA2 context window. As shown in Fig.~\ref{fig:shot_number}, we make two observations: \textbf{(1)} More  CoT shots, facilitated by {CoT-Influx}, indeed boosts LLM math reasoning performance, particularly for larger  LLMs. On LLaMA2-13B, by inputting 48 CoTs, we significantly outperform the standard few-shot CoT and TopK by 4.55\% and 8.72\%, respectively. \textbf{(2)} There is an optimal number of CoT shots for {CoT-Influx}. Its peak performance on LLaMA2 7B and 13B are at 40 and 48 shots, respectively. We attribute this to two potential reasons. First, an extremely large number of shots complicates {CoT-Influx}'s optimization. Second, there may be an upper limit to improving LLM reasoning capability through few-shot learning. 

\textbf{Comparison with state-of-the-art baselines}.  Table~\ref{tab:maintable} and Table~\ref{tab:other_math} present the comparison results of {CoT-Influx} with state-of-the-art baselines across LLaMA2 family and 5 mathematical datasets, highlighting the following observations: \textbf{(1)} by utilizing more few-shot CoTs that are twice the  LLM context window, {CoT-Influx} significantly outperforms all baselines, with 2.13\% to 4.55\% absolute improvements. \textbf{(2)} Despite using fewer input tokens,  {CoT-Influx} consistently outperforms retrieval baselines by 1.36\% to 14.09\% absolute improvements. This is because our compressed tokens indicate more informational CoT examples without redundancy. In contrast, they select entire examples, which may contain redundant tokens, based on semantic similarity between the target question and CoT examples, without considering the different CoT preference of the target LLM. \textbf{(3)} {CoT-Influx} significantly outperforms prompt compression baselines in preserving the most crucial tokens for math reasoning, while methods like Selective Context and LLMLingua suffer accuracy declines due to difficulties in maintaining few-shot prompt structure. GPT-4 tends to prune essential reasoning steps, which negatively impacts CoT effectiveness. 

We further demonstrate the effectiveness of {CoT-Influx} by comparing LLaMA2-70B with larger size LLMs on GSM8K. As shown in Table~\ref{tab:compare_llm}, {CoT-Influx} significantly boosts LLM reasoning capabilities. Remarkably, without any fine-tuning, LLaMA2-70B with {CoT-Influx}  outperform much larger LLMs.  LLaMA2-70B surpasses GPT-3.5 with an absolute improvement of 2.5\%.

 \textbf{Compatible with existing reasoning prompts.}  As a method to improve LLM reasoning capability, {CoT-Influx} is complementary with other advanced reasoning-based prompts. To prove this, we apply self-consistency~\cite{selfconsistency} and self-verification~\cite{weng2023large} to compressed prompts generated by CoT-influx. For evaluation efficiency, we sampled 20 times. As  Table~\ref{tab:majorvote} shows, applying self-consistency and self-verification further improve LLaMA2's performance on GSM8k. 

\subsection{Ablation Study and Analysis}

\textbf{The effectiveness of MRD$^3$ dataset}.  Beyond our pruner, we introduce   MRD$^3$ dataset, which is evolved by GPT-4 for diverse reasoning steps and difficulties. We compare with two baselines: (1) MRD$^3$ without evolution, excluding GPT-4 evolved examples, and (2) the human-labeled GSM8K training set, which excludes  GPT-4's reformatted generation. We apply our pruner on these datasets under the same setting. As shown in  Table~\ref{tab:data_ablation},  both GPT-4 generated and evolved CoT examples are vital for improving the reasoning performance.

\textbf{Ablation study on coarse-to-fine pruner}. Our pruner operates at both shot and token levels to fully exploit redundancy within CoT examples. To verify the effectiveness, we conduct experiments with only shot or token pruner under the same setting. As shown in Table~\ref{tab:pruner_ablation}, removing any pruning stage decreases performance. Notably, removing token-only pruning causes a larger accuracy drop than shot-only pruning, indicating that shot-level redundancy is easier for the pruner to learn. 

\textbf{Token pruning ratios}. We now investigate token pruning ratios by our pruner. Fig.~\ref{fig:pruning_ratio} shows the remaining token length for LLaMA2-70B after our pruner. In total, we achieve a 4.28$\times$ pruning ratio, with shot pruner contributing a 3.87$\times$ ratio. The results  suggest that our pruner favors pruning more coarse-grained shots over fine-grained tokens. 

\textbf{Inference cost}. {CoT-Influx} is a lightweight plug-and-play module, including a 336MB BERT-Large model and a tiny 4MB coarse-to-fine pruner. We measure its additional inference cost.  
Table~\ref{tab:inference_cost} shows the total inference latency and GPU memory required to run LLaMA2-7B with different methods on GSM8K, measured on a single NVIDIA A100 GPU. The results reveal that {CoT-Influx} introduces a negligible 1.4GB additional memory and a 1.7\% increase in latency. This is more effective than prompt compression baselines, such as Selective Context and LLMLingua, which require significantly higher latency and more GPU memory, potentially hindering efficient deployment.

\textbf{Implications}. Our analysis of retained CoT examples and tokens yields the following insights: \textbf{(1)} More capable LLMs favor harder CoT examples, while smaller LLMs opt for  simpler ones. \textbf{(2)} Numerical and format tokens are essential for math reasoning. Function words like \textit{with}, \textit{the}, \textit{then}, and irrelevant background context such as \textit{theater} can be pruned without affecting reasoning capability. 

\section{Conclusion}

We present {CoT-Influx}, a plug-and-play module that improves LLM math reasoning by pruning unnecessary few-shot examples at  shot and token levels for a more effective input context.  To train the module, we use reinforcement learning to optimize a math reasoning-specific reward with GPT-4 evolved CoT dataset MRD$^3$. Extensive experiments on various datasets and LLMs compared with state-of-the-art baselines demonstrate the effectiveness of our method. This paper highlights the vast potential of few-shot CoT prompting in augmenting LLMs' math reasoning abilities.  

\section*{Limitations}
As in-context learning with LLM heavily relies on the selected examples in the prompt, the performance of {CoT-Influx} can be influenced by the quality of CoT generation. Despite this, {CoT-Influx} still demonstrates strong performance on our GPT4-evolved dataset MRD$^3$. We currently use BERT to obtain the feature embedding of a CoT example, which cannot handle long-sequence examples exceeding 512 tokens. We will take these limitations into account and mitigate them in future work. 

{

\end{document}
